\section{Derivations}\label{app:derivations}

\subsection{Freeman--Tukey residuals}\label{app:ft}

Let $n\sim\mathrm{Poisson}(E)$. Freeman and Tukey \cite{FreemanTukey1950} showed that $\sqrt{n}+\sqrt{n+1}$ has variance $1+O(E^{-1})$ and mean $\sqrt{4E+1}+O(E^{-3/2})$, so the residual $r=\sqrt{n}+\sqrt{n+1}-\sqrt{4E+1}$ of \cref{eq:ft} is approximately standard normal for moderate $E$ and remains well behaved, with mean and variance close to their limits, down to $E$ of order one. For a cell with $E\to0$ the residual tends to $r=\sqrt{n}+\sqrt{n+1}-1$, which is zero with probability $1-O(E)$; such cells contribute little to either statistic unless the attack populates them, which is exactly the situation for, for example, double clicks generated by back-scattered light. Multinomial cells with small cell probability behave like Poisson cells; for the no-click cells with probability close to one we do not form residuals, because they are determined by the click cells of the same setting.

\subsection{Variance inflation from estimated bias}\label{app:inflation}

Let $r_\tau=\mu+\varepsilon_\tau$ with independent $\varepsilon_\tau$ of unit variance per component, and let $b=W^{-1}\sum_{w=1}^W r_w$ be estimated on $W$ commissioning windows independent of the $M$ monitoring windows. Then
\begin{equation}
\mathrm{Var}\Bigl(\sum_{\tau=1}^M (r_\tau-b)\Bigr)=M+\frac{M^2}{W}=M\Bigl(1+\frac{M}{W}\Bigr)
\end{equation}
per component, which is the normalisation in \cref{eq:shift,eq:scorestat}. The moving-reference difference of \cref{eq:moving} has, by the same argument with the reference block of length $R$ in place of the commissioning windows, variance $M(1+M/R)$.

\subsection{Fisher information of the multinomial outcome model}\label{app:fisher}

For one setting with $N_s$ pulses and outcome probabilities $p_o(\vartheta)$, the score is $\sum_o n_o\nabla\log p_o$. Since $\sum_o p_o=1$ for all $\vartheta$, $\sum_o p_o\nabla\log p_o=\nabla\sum_o p_o=0$, so the score equals $\sum_o (n_o-N_sp_o)\nabla\log p_o$, which is the form \eqref{eq:score}. Its covariance under the model is
\begin{equation}
\mathrm{Cov}\Bigl(\sum_o n_o g_o\Bigr)=N_s\Bigl[\sum_o p_og_og_o\tr-\Bigl(\sum_o p_og_o\Bigr)\Bigl(\sum_o p_og_o\Bigr)\tr\Bigr]=N_s\sum_o p_og_og_o\tr,
\end{equation}
with $g_o=\nabla\log p_o$, because the second term vanishes by the same identity. Summing over settings gives the information $I_t$ of \cref{sec:score}. The derivatives are evaluated by central finite differences with step $10^{-4}$ in the transformed coordinates of \cref{eq:directions}.

\subsection{Efficient score and the nuisance parameter}\label{app:efficient}

Partition the score as $(u_\psi,u_\eta)$ and the information accordingly. The efficient score $u_\psi-I_{\psi\eta}I_{\eta\eta}^{-1}u_\eta$ is the residual of the $L^2$ projection of $u_\psi$ onto $u_\eta$ and is therefore uncorrelated with $u_\eta$. A first-order error $\delta$ in the nuisance changes $u_\psi$ by $-I_{\psi\eta}\delta$ and $u_\eta$ by $-I_{\eta\eta}\delta$ to first order, so the efficient score is unchanged to first order; this is why profiling $\eta$ per window does not bias the statistic. The covariance of the efficient score is the Schur complement $I_{\psi\psi}-I_{\psi\eta}I_{\eta\eta}^{-1}I_{\eta\psi}$ \cite{Neyman1959,LehmannRomano2005}.

\section{Twin and attack parameters}\label{app:params}

\begin{table}[h]
\centering
\caption{Sampling ranges of the twin. Ranges are uniform unless marked log-uniform. Effects marked $\dagger$ are not represented in the honest model of the monitor.}
\label{tab:twin}
\small
\begin{tabular}{@{}>{\raggedright\arraybackslash}p{5.0cm}>{\raggedright\arraybackslash}p{5.6cm}>{\raggedright\arraybackslash}p{4.4cm}@{}}
\toprule
Quantity & Range & Note \\
\midrule
Signal mean photon number $\mu_\mathrm{s}$ & 0.40--0.60 & \\
Decoy ratio $\mu_\mathrm{d}/\mu_\mathrm{s}$ & 0.15--0.30 & vacuum $\mu_\mathrm{v}=0$ \\
Intensity probabilities $(p_\mathrm{s},p_\mathrm{d},p_\mathrm{v})$ & (0.70, 0.20, 0.10) & bases and bits uniform \\
Fibre length & 5--60\,km (long links: 60--120\,km) & attenuation 0.18--0.22\,dB/km \\
Receiver efficiency $\eta_B$ & 0.10--0.50 (long links: 0.05--0.20) & \\
Optical misalignment $e_Z,e_X$ & 0.5--2\,\%, basis factor 0.8--1.25 & \\
Dark-count probability $p_{dj}$ & $2\times10^{-7}$--$5\times10^{-6}$ per gate & log-uniform, per detector \\
Relative efficiency $\kappa$ & 0.85--1.15 & \\
Afterpulse probability $p_{\rm ap}$ & 0--3\,\% & datasheet value known to $\pm10$\,\% \\
Intensity jitter$^\dagger$ & 0.2--1\,\% per window & independent per intensity \\
Transmittance random walk & innovation 0.1--0.6\,\% per window & AR(1), coefficient 0.995 \\
Periodic transmittance modulation & amplitude 0--4\,\%, period 150--600 windows & \\
Polarisation drift$^\dagger$ & innovation 0--$2\times10^{-4}$ per window & AR(1), coefficient 0.98 \\
Dark-count temperature modulation$^\dagger$ & amplitude 0--10\,\% & same period as above \\
Pulses per window $N$ & $10^7$ & 1\,s at 10\,MHz \\
Windows per link & 200, of which 60 commissioning & \\
\bottomrule
\end{tabular}
\end{table}

\begin{table}[h]
\centering
\caption{Attack strength ranges (log-uniform unless stated). Onsets are uniform between 20 and 100 windows after commissioning.}
\label{tab:attackranges}
\small
\begin{tabular}{@{}lll@{}}
\toprule
Attack & Strength parameter & Range \\
\midrule
Intercept--resend & intercepted fraction $f$ & 0.002--0.2 \\
Photon-number splitting & attacked fraction $f$ & 0.01--1 \\
Detector blinding & blinded fraction $f$ & 0.01--1 \\
Time-shift & shifted fraction $f$; mismatch $r$ & 0.02--1; $r$ uniform 0.1--0.5 \\
Trojan horse & back-scatter probability $q_{\rm bs}$ per gate & $10^{-8}$--$3\times10^{-6}$ \\
RNG bias & amplitude $\delta$; period $T$; phase $\phi$ & $5\times10^{-5}$--$2\times10^{-3}$; $T$ uniform 10--80; $\phi$ uniform \\
\bottomrule
\end{tabular}
\end{table}

\section{Supplementary simulation results}\label{app:results}

\begin{table}[h]
\centering
\caption{Window-level AUROC by attack group (attack-active windows against windows of attack-free test links). Mean\,$\pm$\,95\,\% confidence interval over held-out seeds.}
\label{tab:auroc}
\small
\begin{tabular}{lcccc}
\toprule
Method & Known & Zero-day & Composite & Long links \\
\midrule
\textbf{Q-TRACE} & \textbf{0.859\,$\pm$\,0.033} & \textbf{0.636\,$\pm$\,0.032} & \textbf{0.858\,$\pm$\,0.034} & \textbf{0.699\,$\pm$\,0.032} \\
XGBoost (sup.) & 0.848\,$\pm$\,0.026 & 0.499\,$\pm$\,0.043 & 0.771\,$\pm$\,0.041 & 0.718\,$\pm$\,0.041 \\
QBER CUSUM & 0.664\,$\pm$\,0.043 & 0.533\,$\pm$\,0.045 & 0.664\,$\pm$\,0.034 & 0.624\,$\pm$\,0.019 \\
Decoy CUSUM & 0.746\,$\pm$\,0.022 & 0.546\,$\pm$\,0.031 & 0.724\,$\pm$\,0.067 & 0.660\,$\pm$\,0.036 \\
QBER static & 0.567\,$\pm$\,0.064 & 0.487\,$\pm$\,0.072 & 0.545\,$\pm$\,0.065 & 0.571\,$\pm$\,0.075 \\
Autoencoder & 0.800\,$\pm$\,0.036 & 0.573\,$\pm$\,0.017 & 0.766\,$\pm$\,0.020 & 0.663\,$\pm$\,0.036 \\
Isolation forest & 0.767\,$\pm$\,0.033 & 0.554\,$\pm$\,0.014 & 0.745\,$\pm$\,0.049 & 0.677\,$\pm$\,0.033 \\
One-class SVM & 0.814\,$\pm$\,0.021 & 0.571\,$\pm$\,0.027 & 0.775\,$\pm$\,0.029 & 0.677\,$\pm$\,0.026 \\
QK one-class SVM & 0.523\,$\pm$\,0.018 & 0.498\,$\pm$\,0.013 & 0.521\,$\pm$\,0.014 & 0.501\,$\pm$\,0.025 \\
\bottomrule
\end{tabular}

\end{table}

\begin{table}[h]
\centering
\caption{As \cref{tab:detection}, at a per-window false-alarm rate of $10^{-2}$.}
\label{tab:detection2}
\small
\resizebox{\textwidth}{!}{\begin{tabular}{lccccccc}
\toprule
Method & Known & Zero-day & Composite & Long links & FA honest & FA loss & FA long \\
\midrule
\textbf{Q-TRACE} & \textbf{73.7\,$\pm$\,7.8} & \textbf{42.3\,$\pm$\,4.0} & \textbf{74.7\,$\pm$\,5.9} & \textbf{49.3\,$\pm$\,9.6} & \textbf{33.2\,$\pm$\,20.5} & \textbf{16.8\,$\pm$\,11.3} & \textbf{20.8\,$\pm$\,18.4} \\
XGBoost (sup.) & 61.7\,$\pm$\,12.5 & 3.0\,$\pm$\,3.7 & 50.0\,$\pm$\,8.4 & 31.6\,$\pm$\,12.4 & 4.4\,$\pm$\,3.2 & 29.6\,$\pm$\,11.4 & 0.8\,$\pm$\,2.2 \\
QBER CUSUM & 31.2\,$\pm$\,4.0 & 4.0\,$\pm$\,3.3 & 23.7\,$\pm$\,6.3 & 12.9\,$\pm$\,4.9 & 4.8\,$\pm$\,5.7 & 4.8\,$\pm$\,5.4 & 0.0\,$\pm$\,0.0 \\
Decoy CUSUM & 5.0\,$\pm$\,1.0 & 0.0\,$\pm$\,0.0 & 2.0\,$\pm$\,2.3 & 0.9\,$\pm$\,1.5 & 0.0\,$\pm$\,0.0 & 17.6\,$\pm$\,5.7 & 0.0\,$\pm$\,0.0 \\
QBER static & 15.4\,$\pm$\,4.8 & 3.6\,$\pm$\,3.3 & 15.7\,$\pm$\,9.3 & 56.0\,$\pm$\,3.6 & 5.6\,$\pm$\,2.1 & 3.2\,$\pm$\,6.5 & 48.0\,$\pm$\,14.5 \\
Autoencoder & 6.3\,$\pm$\,3.6 & 0.0\,$\pm$\,0.0 & 4.0\,$\pm$\,3.1 & 1.8\,$\pm$\,3.0 & 0.0\,$\pm$\,0.0 & 20.8\,$\pm$\,16.2 & 0.0\,$\pm$\,0.0 \\
Isolation forest & 7.0\,$\pm$\,6.6 & 1.1\,$\pm$\,1.9 & 7.3\,$\pm$\,10.4 & 0.9\,$\pm$\,1.5 & 0.8\,$\pm$\,2.2 & 40.0\,$\pm$\,12.7 & 0.0\,$\pm$\,0.0 \\
One-class SVM & 0.0\,$\pm$\,0.0 & 0.0\,$\pm$\,0.0 & 0.0\,$\pm$\,0.0 & 0.0\,$\pm$\,0.0 & 0.0\,$\pm$\,0.0 & 0.0\,$\pm$\,0.0 & 0.0\,$\pm$\,0.0 \\
QK one-class SVM & 47.6\,$\pm$\,11.4 & 50.9\,$\pm$\,4.1 & 44.7\,$\pm$\,9.8 & 50.7\,$\pm$\,11.3 & 73.2\,$\pm$\,5.7 & 72.8\,$\pm$\,12.4 & 71.2\,$\pm$\,5.4 \\
\bottomrule
\end{tabular}
}
\end{table}

\begin{table}[h]
\centering
\caption{Test-split attacks at a per-window false-alarm rate of $10^{-3}$: percentage of runs detected, median delay of detected runs, and mean key generated between onset and first alarm (or end of run).}
\label{tab:delaykey}
\small
\begin{tabular}{lccc}
\toprule
Method & Runs detected (\%) & Median delay (windows) & Key exposed (Mbit) \\
\midrule
\textbf{Q-TRACE} & \textbf{50.1} & \textbf{0} & \textbf{6.10} \\
XGBoost (sup.) & 29.6 & 5 & 9.70 \\
QBER CUSUM & 15.8 & 14 & 11.07 \\
Decoy CUSUM & 1.0 & 74 & 12.13 \\
QBER static & 8.4 & 0 & 11.70 \\
Autoencoder & 1.5 & 0 & 12.10 \\
Isolation forest & 2.7 & 38 & 11.88 \\
One-class SVM & 0.0 & -- & 12.14 \\
QK one-class SVM & 5.8 & 44 & 11.81 \\
\bottomrule
\end{tabular}

\end{table}

The key exposed is the key that the decoy-state analysis of \cref{eq:rate} would certify between the onset of an attack and the first alarm. It weights early detection of attacks on high-rate links, and it is the quantity that the security of a deployed system depends on when alarms are used to discard blocks.

\section{Hardware details}\label{app:hardware}

\begin{table}[h]
\centering
\caption{Calibration of the \HwLinks{} links on \HwBackend{} at the time of the run (backend-reported values).}
\label{tab:hwcal}
\small
\begin{tabular}{lccc}
\toprule
Quantity & Median & Min & Max \\
\midrule
CZ error (link pair) (\%) & 0.26 & 0.12 & 0.60 \\
Readout error, data qubit (\%) & 0.74 & 0.29 & 3.31 \\
Readout error, ancilla (\%) & 2.00 & 0.59 & 4.77 \\
$T_1$, data qubit ($\mu$s) & 164 & 68 & 364 \\
$T_2$, data qubit ($\mu$s) & 76 & 18 & 375 \\
$T_1$, ancilla ($\mu$s) & 170 & 98 & 440 \\
$T_2$, ancilla ($\mu$s) & 113 & 26 & 366 \\
\bottomrule
\end{tabular}

\end{table}

The commissioning QBER of the links ranged from \HwQberCalMin\,\% to \HwQberCalMax\,\% with a median of \HwQberCalMed\,\%. The highest value belongs to a link whose readout had degraded since the backend's last calibration; it was retained, because a monitor calibrated per link must cope with such links.

\begin{table}[h]
\centering
\caption{Pre-registered analysis (reference: commissioning passes). Fraction of link-epochs flagged at $\alpha=10^{-3}$. Every test flags most honest link-epochs; see \cref{sec:hw-pre}.}
\label{tab:hwpre}
\small
\begin{tabular}{lccccc}
\toprule
Probe & $n$ & $\Delta Q_{\rm meas}$ (\%) & Q-TRACE & QBER/basis & QBER pooled \\
\midrule
Honest (no probe) & 66 & 1.09 & 0.92 & 0.95 & 0.92 \\
\midrule
$\theta=0.15$ & 14 & 1.48 & 1.00 & 1.00 & 1.00 \\
$\theta=0.25$ & 13 & 2.06 & 1.00 & 0.92 & 1.00 \\
$\theta=0.40$ & 14 & 3.10 & 1.00 & 1.00 & 1.00 \\
$\theta=0.60$ & 14 & 5.53 & 1.00 & 1.00 & 1.00 \\
$\theta=0.90$ & 13 & 10.13 & 1.00 & 1.00 & 1.00 \\
$\theta=1.57$ & 13 & 24.17 & 1.00 & 1.00 & 1.00 \\
\midrule
$\gamma=0.03$ & 13 & 4.06 & 1.00 & 1.00 & 1.00 \\
$\gamma=0.06$ & 12 & 4.47 & 1.00 & 1.00 & 1.00 \\
$\gamma=0.12$ & 12 & 7.15 & 1.00 & 1.00 & 1.00 \\
$\gamma=0.25$ & 14 & 11.06 & 1.00 & 1.00 & 1.00 \\
\bottomrule
\end{tabular}

\end{table}

\begin{table}[h]
\centering
\caption{Condition-matched reference: shots per circuit at which half of the attacked link-epochs are flagged at $\alpha=10^{-3}$ (left; 50 is the smallest value tested), and information gathered by Eve before that point, in bits of sifted key (right).}
\label{tab:hwleak}
\small
\begin{tabular}{lcccccc}
\toprule
 & \multicolumn{3}{c}{Shots per circuit to detection} & \multicolumn{3}{c}{Bits leaked before alarm} \\
\cmidrule(lr){2-4}\cmidrule(lr){5-7}
Probe & Q-TRACE & QBER/basis & QBER pooled & Q-TRACE & QBER/basis & QBER pooled \\
\midrule
$\theta=0.15$ & 511 & 432 & 2851 & 14.9 & 12.6 & 83.0 \\
$\theta=0.25$ & 340 & 369 & 694 & 19.7 & 21.3 & 40.1 \\
$\theta=0.40$ & 73 & 69 & 88 & 13.2 & 12.5 & 15.8 \\
$\theta=0.60$ & 50 & 50 & 50 & 18.2 & 18.2 & 18.2 \\
$\theta=0.90$ & 50 & 50 & 50 & 39.0 & 39.0 & 39.0 \\
$\theta=1.57$ & 50 & 50 & 50 & 69.9 & 69.9 & 69.9 \\
\midrule
$\gamma=0.03$ & 60 & 50 & 76 & 0.4 & 0.4 & 0.5 \\
$\gamma=0.06$ & 50 & 50 & 50 & 1.1 & 1.1 & 1.1 \\
$\gamma=0.12$ & 50 & 50 & 50 & 1.9 & 1.9 & 1.9 \\
$\gamma=0.25$ & 50 & 50 & 50 & 8.3 & 8.3 & 8.3 \\
\bottomrule
\end{tabular}

\end{table}

\begin{table}[h]
\centering
\caption{Eve's information measured on her ancilla. $I_{\rm match}$: mutual information per sifted round in which her probe basis matched Alice's basis (Miller--Madow corrected \cite{Miller1955}); $P_{\rm guess\,err}$: her guessing error on those rounds; $I_{\rm sift}$: information per sifted round averaged over all sifted rounds.}
\label{tab:hwinfo}
\small
\begin{tabular}{lcccc}
\toprule
Probe & $I_{\rm match}$ (meas.) & $I_{\rm match}$ (ideal) & $P_{\rm guess\,err}$ & $I_{\rm sift}$ (meas.) \\
\midrule
$\theta=0.15$ & 0.015 & 0.016 & 0.424 & 0.007 \\
$\theta=0.25$ & 0.029 & 0.045 & 0.399 & 0.014 \\
$\theta=0.40$ & 0.090 & 0.112 & 0.324 & 0.045 \\
$\theta=0.60$ & 0.182 & 0.244 & 0.255 & 0.091 \\
$\theta=0.90$ & 0.390 & 0.505 & 0.148 & 0.195 \\
$\theta=1.57$ & 0.699 & 1.000 & 0.053 & 0.349 \\
\midrule
$\gamma=0.03$ & 0.004 & 0.015 & 0.484 & 0.002 \\
$\gamma=0.06$ & 0.011 & 0.031 & 0.471 & 0.006 \\
$\gamma=0.12$ & 0.019 & 0.063 & 0.450 & 0.010 \\
$\gamma=0.25$ & 0.083 & 0.138 & 0.386 & 0.041 \\
\bottomrule
\end{tabular}

\end{table}

\section{Reproducibility}\label{app:repro}

All results were produced with Python 3.12, NumPy 2.5, SciPy 1.18, pandas 3.0, scikit-learn 1.9, XGBoost 3.4, Qiskit 2.5.2, Qiskit Aer 0.17.2 and Qiskit IBM Runtime 0.49.0. The repository provides \texttt{make benchmark} (all held-out seeds), \texttt{make rehearsal} (noisy simulation of the hardware experiment), \texttt{make figures} (all tables and figures) and the notebook \texttt{notebooks/01\_qpu\_hardware\_run.ipynb} used for the hardware run, including its budget control: a pilot job measures the cost per shot and the main job is sized to a user-set limit. Every random draw in the benchmark derives from the seed, so a seed reproduces its result files exactly with the pinned package versions. The numbers in this paper are generated from the result files by the script \texttt{paper/build\_numbers.py}.
